% ===============================================================
\section{Experimental Setup}\label{sec:eval}
% ===============================================================

\subsection{Research Questions}\label{subsec:rqs}

We organise the evaluation around five research questions that separate the
two claims of the paper---that routing among evidence mechanisms is both
learnable and useful, and that a selective second stage adds value beyond
route correction.

\begin{description}
\item[RQ1 (End-to-end quality).] Does adaptive routing improve answer quality
over fixed-strategy and oracle baselines, and does Stage~2 contribute beyond
mere routing correction?
\item[RQ2 (Routing accuracy).] How accurately does the Stage~1 classifier
select routes under cross-validation, compared with rule-based and learned
baselines?
\item[RQ3 (Ambiguity handling).] Can the system avoid unsafe retrieval by
clarifying when the legal target is underspecified?
\item[RQ4 (Threshold ablation).] How do the key thresholds affect routing
accuracy and the Stage~2 trigger rate?
\item[RQ5 (Cost--quality).] Does selective Stage~2 invocation provide a better
trade-off than always-on or always-off verification?
\end{description}

\subsection{Systems and Baselines}\label{subsec:systems}

Table~\ref{tab:systems} defines all compared systems. All end-to-end
experiments run on the \texttt{phapdien\_strict} test set of $600$ queries,
using the single Pháp~điển corpus with one shared retrieval and generation
configuration, so that any observed difference reflects routing decisions
rather than infrastructure. The Oracle Router dispatches every query using its
gold route label under identical retrieval and generation settings; comparing
the oracle against the Two-stage Hybrid therefore quantifies how much answer
quality is attributable to Stage~2 reasoning beyond correct routing.

% NOTE(M4): the "Oracle + Stage 2" and "Always-on Stage 2" rows isolate the
% chain-of-thought contribution and the value of SELECTIVE triggering,
% respectively. Both have now been run (see Table~\ref{tab:end_to_end}).
\begin{table}[!htbp]
\centering
\caption{Systems compared in this work}\label{tab:systems}
\begin{tabularx}{\columnwidth}{L{2.7cm}Y}
\toprule
\textbf{System} & \textbf{Description} \\
\midrule
Pure Vector & Always dense retrieval from Chroma \\
Pure Graph & Always graph traversal in Neo4j (Text-to-Cypher) \\
Pure Hybrid & Always combines dense and graph evidence \\
Single-stage Router & XGBoost Stage~1 only; no LLM verifier \\
Two-stage Hybrid & Stage~1 + selective Stage~2 (full system) \\
Always-on Stage~2 & Stage~1 + unconditional Stage~2 verification (trigger-policy ablation) \\
Oracle Router & Gold route labels; retrieval only, no Stage~2 \\
Oracle + Stage~2 & Gold route labels \emph{with} Stage~2 reasoning (CoT-isolation ablation) \\
\bottomrule
\end{tabularx}
\end{table}

\subsection{Evaluation Metrics}\label{subsec:metrics}

\paragraph{Answer quality.}
For answer quality we report Exact Match (EM) and token-level F1 computed by a
single \texttt{evaluate\_prediction(gold\_context, answer)} routine shared by
all systems; by convention F1$\,=0$ when precision and recall are both zero.
Because token-F1 against a gold context rewards lexical overlap and can
penalise a correct answer phrased differently from the reference---a bias that
disadvantages systems whose evidence is assembled from structured graph
context rather than copied from a passage---we additionally report a semantic
metric to complement lexical F1.
% TODO(M6): choose ONE semantic metric and report it alongside EM/F1 for every
% system in Table~\ref{tab:end_to_end}. Options: (i) BERTScore-F1 with a
% Vietnamese encoder (bib entry zhang2020bertscore is already in biblio.bib but
% currently uncited -- citing it here makes it used); or (ii) an LLM-as-judge
% correctness score with a fixed rubric. State the model and settings used.
We report BERTScore-F1~\cite{zhang2020bertscore} with a Vietnamese encoder as
this semantic complement. Hit@$k$ measures retrieval coverage (the gold source
appears within the top $k$ results) over the source-labelled queries. Because
the dense index, the graph store, and the gold annotations historically used
four incompatible identifier formats, both retrieved and gold references are
first mapped to a canonical \texttt{law::\allowbreak document::\allowbreak
article::\allowbreak number} form by a deterministic normaliser before
matching (strict mode: document code \emph{and} article number must agree);
references the normaliser cannot resolve are counted as misses and flagged,
rather than silently matched (Appendix~\ref{app:repro}).

\paragraph{Route classification.}
For route classification we report accuracy, per-class F1, and Macro-F1 over
the class set $\mathcal{R}=\{d_v,d_g,d_h\}$; Macro-F1 is emphasised because the
dataset is dominated by the dense class~\cite{powers2011evaluation}. To avoid
reporting an over-optimistic single-split figure, the headline routing result
is obtained by \emph{stratified five-fold cross-validation}, and we report the
mean and standard deviation across folds.

\paragraph{Clarification.}
For clarification we report precision, recall, and F1 of the \texttt{clarify}
decision on the ambiguity benchmarks, computed against the binary
ambiguous/answerable gold label.

\paragraph{Cost.}
For cost we report average and 95th-percentile (p95) latency, the Stage~2
trigger rate, and the graph-to-dense fallback rate. To attribute latency to its
sources rather than to the pipeline as a whole, we additionally decompose
end-to-end latency into a routing component (Stage~1, optional Stage~2) and a
retrieval-plus-generation component (Section~\ref{subsec:rq5}).

\paragraph{Statistical reporting.}
Because all end-to-end scores are means over a finite test set ($N=600$),
point estimates alone can overstate the separation between systems. We
therefore accompany the headline F1 and Hit@$1$ figures with $95\%$
confidence intervals obtained by bootstrap resampling of the per-query scores
($B=1000$ resamples, fixed seed~$42$), and we regard two systems as
distinguishable only when their intervals do not overlap.
